\subsection*{Significance statement}
Understanding how species diversify or pathogens spread requires linking phylogenetic trees to predictors such as biological traits and environmental features.
We present an unsupervised Bayesian neural network that infers the complex links between predictors and evolutionary parameters---such as speciation, extinction, and transmission rates---directly from sequence data.
The model does not require training data, enabling inferences otherwise intractable.
Our framework outperforms state-of-the-art alternatives, while remaining interpretable through explainable artificial intelligence.
Empirical applications in epidemiology reveal nonlinear effects of travel among countries on the early spread of SARS-CoV-2 in Europe, while a macroevolutionary application shows selective extinction events in the diversification history of New World monkeys.
This work expands the reach of phylodynamics in biology and public health.
